% ============================== CHAPTER 1 ===================================
\chapter{Introduction}
\label{ch:introduction}

This chapter establishes the setting in which \setu{} was built. It describes the
linguistic and infrastructural problem that motivates an offline Indic
translator, summarises the state of the art the project builds upon, states the
objectives that the system was measured against, and closes with a guide to the
remaining chapters of this report.

\section{Overview}
\label{sec:overview}

The Eighth Schedule of the Constitution of India recognises twenty-two
languages. Together with several hundred dialects, these languages define the
everyday linguistic reality of more than a billion people. Digital services,
however, remain overwhelmingly English-first. The gap between the language a
citizen speaks and the language a service responds in is not a matter of
convenience; it decides whether a person can read a government form, understand
a prescription, or follow an advisory during a disaster.

Machine translation (MT) is the standard technical answer to that gap, and for
Indian languages the answer has become a very good one. IndicTrans2
\cite{ref:indictrans2} covers all twenty-two scheduled languages at a quality
that was unthinkable five years ago, trained on the Bharat Parallel Corpus
Collection (BPCC) and evaluated on the IN22 and FLORES-200 benchmarks. Earlier
multilingual efforts (mT5-based Indic systems \cite{ref:mt5indic}, Indic-to-Indic
pivoting systems \cite{ref:indic2indic}, and one-to-many low-resource systems for
Assamese and Bengali \cite{ref:assamese}) established that shared multilingual
representations transfer usefully across related Indian languages.

Every one of those systems, however, shares an architectural assumption that
this project rejects: that the model runs in a data centre and the user's text
travels to it. Two consequences follow. The first is availability. A cloud
translator is useless exactly where translation matters most: the rural
primary health centre with intermittent connectivity, the border post beyond
cellular coverage, the relief camp after a cyclone has taken the towers down.
The second is privacy. A cloud translator sees every sentence a user submits.
For a citizen dictating a medical complaint or a legal grievance, that is a
meaningful disclosure, and it is a disclosure with no technical necessity.

The natural alternative is to run the model on the device. Here the arithmetic
is unforgiving. The flagship IndicTrans2 checkpoints hold roughly a billion
parameters and occupy several gigabytes in 32-bit floating point; even the
distilled 200-million-parameter variants need close to a gigabyte of resident
memory. A mid-range Android handset or an ARM Cortex-A55 single-board computer
cannot host either. Closing the gap therefore requires a model one order of
magnitude smaller that nevertheless retains most of the teacher's translation
quality.

Model compression research offers several routes to such a model.
Knowledge distillation trains a compact student to imitate a larger teacher;
DistilBERT \cite{ref:distilbert} demonstrated the approach for encoders, and for
translation specifically, sequence-level knowledge distillation replaces human
reference targets with the teacher's own decoded output. SMaLL-100
\cite{ref:small100} showed that a deliberately shallow multilingual translator can
remain competitive for low-resource languages. Work on neural translation for
mobile devices \cite{ref:mobilenmt} and on shallow-to-deep training
\cite{ref:shallow2deep} addressed the latency and stability problems that arise
when translation models are shrunk. In parallel, Direct Preference Optimization
(DPO) \cite{ref:dpo} made preference-based training practical without a reward
model or reinforcement learning, and subsequent work applied it to neural
machine translation with minimum Bayes risk decoding \cite{ref:dpomt}.

\setu{} (from the Sanskrit \dev{सेतु}, meaning \emph{bridge}) sits at the
intersection of these lines of work. It takes IndicTrans2 as a teacher and
distils it into compact students of 52.6 million parameters, one per translation
direction, quantised to INT4 ONNX and executed by ONNX Runtime with a local
SentencePiece tokenizer. Nothing leaves the device at inference time.

The project's original thesis was that preference distillation would be the
mechanism: generate several teacher hypotheses per source sentence, rank them by
chrF, pair the best against the rest, and train the student with DPO. That
thesis was tested directly against the standard sequence-level knowledge
distillation baseline at matched student size, and it lost decisively:
17.09 BLEU for sequence-level distillation against 11.10 BLEU for the
preference-trained student. The diagnosis, developed in
Chapter~\ref{ch:results}, is that the limiting factor for a compact Indic
translator is not the training objective but the noise in web-mined human
references. Training on the teacher's clean one-best output removes that noise,
and the student improves accordingly. This report presents that refutation as a
result in its own right, because a negative result obtained under controlled
conditions is more useful to the field than an unexamined assumption.

The delivered system trains each direction on approximately 497{,}000
teacher-distilled sentence pairs, quantises the result to 103.9\,MB, and serves
it through a REST API, an offline progressive web application, a browser client
and a command-line tool. Eleven scheduled languages are covered in both
directions. Because no ungated Indic-to-Indic teacher was available, pairs such
as Hindi--Bengali are served by chaining two students through English, which
turns 22 trained models into 132 usable translation directions.

\section{Objectives}
\label{sec:objectives}

The project set out to design and deliver a multilingual distillation pipeline
that converts a large Indic translation teacher into compact, quantised students
capable of fully offline inference on edge hardware; to determine empirically
whether preference-based distillation or sequence-level knowledge distillation
produces the better compact student at matched size; to cover a substantial
subset of the twenty-two scheduled languages in both directions and extend
coverage to every Indic-to-Indic pair without training a quadratic number of
models; and to expose the resulting engine through interfaces that a
non-specialist can actually use.

Stated as measurable targets, the objectives are:

\begin{enumerate}[label=\textbf{O\arabic*.}, leftmargin=0.75in]
  \item \textbf{Quality.} Each student shall achieve a BLEU score at least
        80\,\% of its teacher's BLEU on the same held-out evaluation set, i.e.\
        $\text{BLEU}_{\text{student}} / \text{BLEU}_{\text{teacher}} \geq 0.80$.
  \item \textbf{Latency.} Translation of a single sentence shall complete in
        under 500\,ms on ARM-class CPU hardware, with no GPU and no accelerator.
  \item \textbf{Size.} The deployable quantised artifact for one direction shall
        not exceed 200\,MB.
  \item \textbf{Offline operation.} Inference shall succeed with all network
        sockets disabled; no user text shall leave the device.
  \item \textbf{Empirical finding.} The preference-distillation hypothesis shall
        be tested against a sequence-level knowledge distillation baseline at
        matched student size, on identical data, with evaluation on real human
        references, and the outcome reported whichever way it falls.
  \item \textbf{Coverage and extensibility.} A substantial subset of the
        scheduled languages shall be supported bidirectionally, all
        Indic-to-Indic pairs shall be reachable, and adding a new language shall
        be a configuration and data change rather than a code change.
\end{enumerate}

Objectives O2, O3 and O4 are met by every trained direction without exception.
Objective O1 is met by 14 of 22 directions; the remaining eight fall between
0.755 and 0.787, and Chapter~\ref{ch:results} analyses why. Objective O5
produced a refutation of the project's own founding hypothesis. Objective O6 is
met for eleven languages, with the remaining eleven blocked by corpus
availability rather than by the architecture.

\section{Purpose, Scope and Applicability}
\label{sec:psa}

\subsection{Purpose}
\label{subsec:purpose}

The purpose of \setu{} is to make high-quality Indic machine translation
available where the network is not, and to do so without requiring the user to
disclose what they are translating.

Three concerns motivate the work. The first is \emph{access}. Connectivity in
India is neither uniform nor reliable, and the populations least well served by
it overlap heavily with the populations least well served in English. A
translator that requires bandwidth systematically fails the users with the
greatest need for it. The second is \emph{privacy}. Translation input is
frequently sensitive: symptoms, grievances, legal matters, personal
correspondence. On-device inference makes the privacy guarantee structural
rather than contractual: there is no server to trust because there is no server.
The third is \emph{sovereignty of infrastructure}. A public institution that
depends on a commercial translation API is exposed to pricing changes, API
deprecations and terms-of-service revisions. A 104\,MB artifact that runs on the
institution's own hardware is not.

Theoretically, the project addresses a question that the compact-model
literature has not settled for translation: when a student is an order of
magnitude smaller than its teacher and the available parallel data is mined
rather than curated, which training signal should the student be given? The
literature offers reference targets, teacher one-best targets and
preference-ranked teacher hypotheses. \setu{} evaluates all three under matched
conditions and reports the outcome.

\subsection{Scope}
\label{subsec:scope}

\textbf{Methodology.} The system distils the ungated rotary IndicTrans2
distilled-200M checkpoints into 52.6-million-parameter Marian sequence-to-sequence
students, one per direction, using sequence-level knowledge distillation on
teacher one-best output. Students are quantised post-training to INT8 and INT4
ONNX and executed by ONNX Runtime. Evaluation uses sacreBLEU (BLEU and chrF)
against real human references held out from the training corpus, with the
student/teacher BLEU ratio as the headline metric.

\textbf{Languages delivered.} Eleven scheduled languages (Hindi, Bengali,
Marathi, Telugu, Tamil, Gujarati, Kannada, Odia, Malayalam, Assamese and Punjabi) each paired bidirectionally with English, giving 22 trained directions. Every
Indic-to-Indic pair among these eleven is served by chaining two students
through English.

\textbf{Assumptions.} The work assumes that the ungated distilled-200M teacher is
an acceptable quality ceiling; that Samanantar, though web-mined and noisy, is
representative enough for distillation; that the student architecture and its
16{,}000-token vocabulary are sized correctly for the data volume available; and
that x86 CPU latency measurements are a reasonable proxy for ARM behaviour.
Each assumption is examined in Chapter~\ref{ch:results}.

\textbf{Limitations, stated up front.} The project is text-to-text only; speech
input and output are out of scope. The remaining eleven scheduled languages are
not delivered, because only Assamese and Punjabi among them have ungated
Samanantar data and the rest require access to the gated BPCC corpus. The
quality ratio is bounded by the ungated 200M teacher, which is itself weaker than
the gated one-billion-parameter flagship. Reported ratios are computed on an
internal Samanantar development slice rather than on FLORES-200, so they are
comparable within this report but not directly comparable to published
IndicTrans2 numbers. ARM latency is projected from x86 measurement rather than
measured on ARM silicon. Finally, INT4 quantisation yields the same on-disk size
as INT8 in ONNX Runtime's CPU execution provider, which lacks a true 4-bit
kernel.

\subsection{Applicability}
\label{subsec:applicability}

\textbf{Direct applications.} In rural healthcare, a doctor or ASHA worker can
converse with a patient across a language boundary on a handset with no
connectivity. In defence and border operations, patrol units operating beyond
cellular coverage can translate without transmitting anything. In education,
multilingual tutoring tools can serve students in their own language without a
per-query cloud cost. In government service delivery, forms, notices and
entitlement documents can be translated at a gram panchayat counter. In disaster
relief, field workers can communicate with affected communities during the exact
window in which the network is down.

\textbf{Indirect contributions.} The distillation pipeline is not
Indic-specific; the same configuration-driven process applies to any
language family with a strong teacher and a weak-reference parallel corpus. The
controlled comparison between preference distillation and sequence-level
distillation contributes a data point to the compact-NMT literature that
practitioners can act on. The English-pivot router demonstrates that $n$
languages can be served in $n^2$ directions with $2n$ models, a result of
practical importance when each additional model costs 104\,MB of device storage.
And the engineering record in this report (particularly the two training
collapses and their diagnosis) documents failure modes that are rarely
published but routinely encountered.

\section{Organisation of the Report}
\label{sec:organisation}

\textbf{Chapter~\ref{ch:literature}, Literature Survey}, reviews twenty papers
published from 2019 onward across mobile and compact NMT, knowledge
distillation, preference optimisation and multilingual Indic translation,
compares the available technologies, and derives the drawbacks of existing
systems, the problem statement and the proposed solution.

\textbf{Chapter~\ref{ch:requirements}, Requirement Engineering}, specifies the
hardware and software used for training and deployment, presents the conceptual
models (use case, sequence, activity, state chart, data flow, entity
relationship and class diagrams) and states the user, system, functional,
non-functional and domain requirements.

\textbf{Chapter~\ref{ch:planning}, Project Planning}, gives the work breakdown
structure, the Gantt schedule across both project phases, a PERT network with
the critical path, the risk register and the division of responsibility within
the team.

\textbf{Chapter~\ref{ch:design}, System Design}, presents the high-level
architecture, decomposes the system into eight modules, specifies the external
interfaces and the core data structures, and gives the algorithms for
distillation, training, quantisation and pivot routing.

\textbf{Chapter~\ref{ch:implementation}, Implementation}, describes the
implementation approach and coding standards, presents the significant code with
commentary, and discusses the optimisations that made training feasible on the
available hardware.

\textbf{Chapter~\ref{ch:testing}, Testing}, sets out the testing strategy and
reports the unit, integration and system test cases with their data and
outcomes.

\textbf{Chapter~\ref{ch:results}, Results, Discussion and Performance
Analysis}, reports the full 22-direction scorecard, the preference-versus-%
sequence-distillation head-to-head, the data-scaling curve, the collapse and
recovery analysis, qualitative translation samples and the user documentation.

\textbf{Chapter~\ref{ch:conclusion}, Conclusion, Applications and Future Work},
summarises the outcome against each objective, lists applications and
limitations, and defines the work that follows.

The report closes with \textbf{References}, an \textbf{Appendix} covering the
publication plan and reproduction instructions, and a \textbf{Glossary}.
